% Teil II: der durchgerechnete Weg. Reihenfolge beibehalten.

\subsection*{Gegeben und gesucht}
Sei $a\in\mathbb R$ fest. Gegeben ist eine auf $\mathbb R$ differenzierbare
Funktion $f$ mit $f'=af$. Gesucht: Herleitung der Gestalt aller solchen
Funktionen über die Hilfsfunktion $g(x)=f(x)\mathrm e^{-ax}$.

\subsection*{Ansatz}
Man bildet $g$, von der man − dank Produktregel und $f'=af$ − zeigen kann,
dass $g'=0$ ist. Nach dem Satz, dass $g'=0$ Konstanz bedeutet, ist $g$
konstant; das stellt man nach $f$ um.

\subsection*{Rechnung}
Die Ableitung von $g(x)=f(x)\mathrm e^{-ax}$ mit der Produktregel:
\begin{equation}
  g'(x)=f'(x)\,\mathrm e^{-ax}+f(x)\,\frac{\mathrm d}{\mathrm dx}\mathrm e^{-ax}
  = f'(x)\,\mathrm e^{-ax}-a f(x)\,\mathrm e^{-ax}.
\end{equation}
Wegen $f'=af$ kürzt sich der zweite Term:
\begin{equation}
  g'(x)=a f(x)\,\mathrm e^{-ax}-a f(x)\,\mathrm e^{-ax}=0.
\end{equation}
Da $g$ differenzierbar und überall die Ableitung null hat, ist $g$ konstant.
Es gibt also ein $\lambda\in\mathbb R$ mit $f(x)\mathrm e^{-ax}=\lambda$.
Umstellen nach $f$ ergibt die gesuchte Gestalt.

\subsection*{a) Ableitung und b) Konstanz}
\begin{ergebnis}
\begin{equation*}
  \frac{\mathrm d}{\mathrm dx}\left(f(x)\mathrm e^{-ax}\right)=0
  \qquad\Rightarrow\qquad
  f(x)\mathrm e^{-ax}\ \text{ist konstant}.
\end{equation*}
\end{ergebnis}

\subsection*{c) Allgemeine Gestalt}
Aus $f(x)\mathrm e^{-ax}=\lambda$ folgt durch Multiplikation mit $\mathrm e^{ax}$:
\begin{ergebnis}
\begin{equation*}
  f(x)=\lambda\,\mathrm e^{ax},\qquad \lambda\in\mathbb R.
\end{equation*}
\end{ergebnis}

\begin{kontrolle}
\textbf{Umkehrung:} Für $f(x)=\lambda\mathrm e^{ax}$ ist $f'(x)=a\lambda\mathrm e^{ax}=af(x)$,
die Bedingung ist also nicht nur nötig, sondern auch hinreichend.
\textbf{Konstante:} Bei $x=0$ ist $g(0)=f(0)\mathrm e^{0}=f(0)$, also
$\lambda=f(0)$; die Konstante ist der Funktionswert im Nullpunkt.
\textbf{Grenzfall:} Für $a=0$ bleibt $f$ konstant, wie es die Formel
$f(x)=\lambda\mathrm e^{0}=\lambda$ liefert.
\end{kontrolle}

\begin{merksatz}
Differentialgleichungen des Typs $f'=af$ löst man über eine
Hilfsfunktion $\mathrm e^{-ax}$, deren Ableitung null man nachweist. Das
Muster „Ableitung null $\Rightarrow$ konstant“ taucht bei allen Charakterisierungen
von Exponentialfunktionen auf.
\end{merksatz}
